% 操作系统原理课内实践项目 · 个人分析（每人半页至一页）
% 版式基线：Paper_Library/计算方法实验报告模板（汕头大学课程报告 LaTeX 模板库）
% 编译：XeLaTeX（中文必须 XeLaTeX）。
\documentclass[12pt,a4paper]{ctexart}
\usepackage{amsmath,amssymb}
\usepackage{graphicx}
\usepackage{booktabs}
\usepackage{array}
\usepackage{geometry}
\usepackage{caption}
\usepackage{float}
\usepackage{enumitem}
\usepackage{hyperref}
\hypersetup{hidelinks}

\geometry{left=2.6cm,right=2.6cm,top=2.6cm,bottom=2.6cm}
\captionsetup{font=small,labelsep=quad}
\renewcommand{\arraystretch}{1.15}
\setlist{nosep}

% 标题（与模板库正文标题同级，不加编号）
\newcommand{\mysection}[1]{\par\vspace{1.0em}\noindent{\large\bfseries #1}\par\vspace{0.45em}}

% 写作提示样式
\newcommand{\hint}[1]{\par\noindent{\small\itshape\kaishu #1}\par}
% 待填写横线
\newcommand{\blank}[1][\linewidth]{\underline{\hspace{#1}}}
\newcommand{\fillline}{\par\vspace{0.55em}\noindent\blank[\linewidth]\par}

\begin{document}

% ==================== 标题 ====================
\begin{center}
    {\zihao{3} \textbf{《操作系统原理》课内实践个人分析}}\\[0.5em]
    {\zihao{4} 题目五：连续内存分配与碎片问题验证}
\end{center}
\vspace{0.2em}
\hrule
\vspace{0.8em}

% ==================== 基本信息 ====================
\begin{table}[H]
    \centering\small
    \begin{tabular}{|c|c|p{3.6cm}|p{5.0cm}|}
        \hline
        \textbf{姓名} & \textbf{学号} & \textbf{组内分工} & \textbf{主要贡献位置（文件/commit）} \\ \hline
        \blank[1.9cm] & \blank[1.9cm] & \blank[3.3cm] & \blank[4.7cm] \\ \hline
    \end{tabular}
\end{table}

% ==================== 一、初始假设 ====================
\mysection{一、我的初始假设}
\hint{我最初怎样判断，为什么这样判断。写讨论之前的第一反应，不必与最终结论一致。}

\hint{示例句式：「我最初认为最佳适应一定比首次适应好，因为它每次挑最合适的块；
但我说不出在什么负载下这个判断会失效。」}
\fillline
\fillline

% ==================== 二、亲自完成的验证 ====================
\mysection{二、我亲自完成的验证}
\hint{我亲自完成了哪项推演或实验，证据在哪里（对应仓库中的具体文件与步骤）。}

\begin{table}[H]
    \centering\small
    \setlength{\tabcolsep}{4pt}
    \begin{tabular}{|>{\raggedright\arraybackslash}p{3.9cm}|>{\raggedright\arraybackslash}p{3.4cm}|>{\raggedright\arraybackslash}p{4.0cm}|>{\raggedright\arraybackslash}p{2.4cm}|}
        \hline
        \textbf{验证项} & \textbf{方法} & \textbf{证据文件位置} & \textbf{结论} \\ \hline
        \blank[3.6cm] & \blank[3.1cm] & \texttt{experiments/}\blank[1.6cm] & \blank[2.1cm] \\ \hline
        \blank[3.6cm] & \blank[3.1cm] & \texttt{experiments/}\blank[1.6cm] & \blank[2.1cm] \\ \hline
        \blank[3.6cm] & \blank[3.1cm] & \texttt{experiments/}\blank[1.6cm] & \blank[2.1cm] \\ \hline
    \end{tabular}
\end{table}

% ==================== 三、认识变化 ====================
\mysection{三、核验后我的认识变化}
\hint{核验后我修正或确认了什么，还有什么不能据此推出。}

\noindent 最初认为：\blank[0.86\linewidth]\fillline
\noindent 证据显示：\blank[0.86\linewidth]\fillline
\noindent 现在认为：\blank[0.86\linewidth]\fillline
\noindent 仍不能据此推出：\blank[0.80\linewidth]

% ==================== 四、自查 ====================
\mysection{四、个人考核自查}

\begin{table}[H]
    \centering\small
    \setlength{\tabcolsep}{4pt}
    \begin{tabular}{|>{\raggedright\arraybackslash}p{4.4cm}|c|>{\raggedright\arraybackslash}p{6.2cm}|c|}
        \hline
        \textbf{评分点} & \textbf{分值} & \textbf{我的对应内容/证据} & \textbf{自评} \\ \hline
        独立解释原理及本人证据 & 15 & \blank[6.0cm] & \blank[1.2cm] \\ \hline
        针对输入变化提出假设或手工推算 & 10 & \blank[6.0cm] & \blank[1.2cm] \\ \hline
        可核对的个人贡献 & 5 & \blank[6.0cm] & \blank[1.2cm] \\ \hline
    \end{tabular}
\end{table}

\hint{「针对输入变化提出假设」可写：把核心案例第 6 步的申请量由 50 改为 30 后，
预测两种策略各分配到哪里、指标如何变化，再运行核对——这是验收时最常被追问的一项。}

\end{document}
